\documentclass[12pt]{amsart}
\usepackage{amsmath,amssymb,amsthm,bm}
\usepackage[usenames,dvipsnames]{xcolor}
\usepackage{graphicx}
%\usepackage{lmodern}
\usepackage[T1]{fontenc}
%\usepackage{textcomp}
\usepackage{pxfonts}
\usepackage{enumerate,verbatim,cite}
\usepackage[margin=1in]{geometry}
%\usepackage{../../tex/pythonhighlight} %https://github.com/olivierverdier/python-latex-highlighting
\usepackage[framed]{mcode}

\usepackage{fancyhdr}
\pagestyle{fancy}
%\addtolength{\headheight}{\baselineskip}
\addtolength{\footskip}{\baselineskip}
\renewcommand{\headrulewidth}{0pt}
\renewcommand{\footrulewidth}{0.4pt}
\fancyhf{}
\fancyfoot[L]{\textit{Last Modified: \today}}
\fancyfoot[C]{\thepage}




\usepackage[pdftex,bookmarks,hyperfigures,colorlinks
						,urlcolor=blue
						,citecolor=blue
						,linkcolor=blue
						,pdfstartview=FitH]{hyperref}


%\usepackage{epsf}
% \topmargin -0.5in \setlength{\textwidth}{6.in}
% \setlength{\textheight}{8.5in}
%\setlength{\evensidemargin}{0.25in}
% \setlength{\oddsidemargin}{0.25in}
\renewcommand{\r}{{\bf r}}
\newcommand{\dery}{\frac{dx}{dt}}
\renewcommand{\k}{{\bf k}}
\newcommand{\be}{\begin{equation}}
\newcommand{\ee}{\end{equation}}
\newcommand{\ns}{\vspace{3mm}}
\newcommand{\Pay}{\text{Payoff}}
\newcommand{\oPay}{\overline{\Pay}}



\title{Applied Math 50 Lab 5: Game Theory
}

\date{}

\begin{document}
\maketitle

\section{Introduction}
In today's computer lab, we will:
\begin{enumerate}
 \item Show how the ideas of game theory that we discussed in class
    can be implemented in a differential equation model, and we will discover
    the analogue of the Nash equilbrium in these models. We will use the models
    to simulate the examples from class.
  \item Study the mathematics of the models in detail and show how to
    calculate the possible Nash equilibria.
  \item Introduce the \textit{AM 50 competition} based on the unique number
    game, and show you how to write functions for it.

\end{enumerate}

\section{A differential equations model of a game}
\noindent In class we discussed game theoretic payoff matrices such as
\[
A=\left( \begin{array}{cc}
a_{11} & a_{12} \\
a_{21} & a_{22} \end{array} \right),
\]
representing the payoffs of player 1 in a two-player game. The value $a_{11}$ is the
payoff if both players choose the first strategy; the value $a_{22}$ is the payoff
if both players choose the second strategy. The value $a_{12}$ is the payoff if
player 1 chooses the first payoff and player 2 chooses the second.   The value $a_{21}$ is the payoff if
player 1 chooses the second payoff and player 2 chooses the first.  Here, we
will consider a symmetric game, where the payoffs for player 2 are given by,
\[
A^T=\left( \begin{array}{cc}
a_{11} & a_{21} \\
a_{12} & a_{22} \end{array} \right).
\]

\noindent We discussed in class Nash's idea that there can be an equilibrium point
of the game, and that in general the equilibrium point is a mixed strategy as
opposed to a pure strategy. Here we will show you an explicit implementation of
this idea in terms of ordinary differential equations.\\

\noindent Let $A$ be the payoff matrix as discussed above. We will imagine that there are
a large number of people (say 1,000) who are all playing the game according to
the payoff matrix. Each of them adopts some mixed strategy described by a
vector $(p,q)$, where $p$ is the probability of playing the first strategy, and
$q$ is the probability of playing the second strategy. Choosing $p=1,q=0$
corresponds to choosing the first strategy, and choosing $p=0,q=1$ is the
second. Anything in-between is a mixed strategy.\footnote{Note that this
framework is slightly more restrictive than that considered in the lecture,
since we assume that all players are playing the same strategy.}\\

\noindent Now let us suppose that if you choose a strategy with a bigger payoff, you are
more likely to choose this strategy in the future, whereas if you choose a
strategy with a smaller payoff, you are less likely to choose this strategy.
This is a sensible model, in that if you find yourself losing a lot of money
playing a game you will change the way that you play it until that doesn't
happen anymore.\\

\noindent A simple model that captures this idea is the following set of ordinary
differential equations. These equations are called the ``replicator dynamics,''
and were introduced by Taylor and Jonker.\footnote{There is a very nice
description of these dynamics in Martin Nowak's book ``Evolutionary dynamics,''
and in his course in the math department.} The model is
\begin{eqnarray}
  \frac{dp}{dt}= p \left( \Pay_1 - \oPay \right) \label{eq:rep1}, \\
  \frac{dq}{dt}=q\left( \Pay_2 - \oPay \right) \label{eq:rep2},
\end{eqnarray}
where $$\Pay_{1}= p a_{11} + q a_{12}$$ and $$\Pay_{2}= p a_{21}+ q a_{22}.$$\\
These are the payoffs expected from playing each of the two strategies. The
quantity $\oPay$ is the {\it average payoff}. This is given by
$$\oPay=p {\Pay}_1+q{\Pay}_2.$$ The model therefore says that if the payoff of
a particular strategy is larger than the average payoff, the probability of
playing this strategy will grow, whereas if the payoff is less than the average
the probability of playing this strategy will shrink. The question we need to
consider is whether these dynamics will lead to an equilibrium, and if it
agrees with the Nash equilibrium discussed in class.




\bigskip
\begin{center}
\colorbox{Goldenrod}{\parbox{15cm}{
{
\textbf{\large Exercise 1}\vspace{1mm} \\
Start with the Matlab program we gave you for solving the SIR model from Lab 1.
Modify these equations to solve the replicator dynamics described above.
Note that in Eqs.~\ref{eq:rep1} \& \ref{eq:rep2} we have that $p+q=1$, because
the total probability of choosing one strategy or the other must sum to unity.
Thus when you choose initial conditions, you need to choose them so that the
initial $p+q=1$. }}}
\end{center}
\bigskip 

\noindent To simulate the equation properly, you need to choose different examples for
the matrix $A$. We propose you start with the three examples we discussed in
class.

%\subsection{The Game of Chicken}
%\noindent Recall for the game of chicken we had the payoff matrix
%\[
%A=\left( \begin{array}{cc}
%-c & b \\
%-b &0 \end{array} \right).
%\]
%
%\bigskip
%\begin{center}
%\colorbox{Goldenrod}{\parbox{15cm}{
%{
%\textbf{\large Exercise 2}\vspace{1mm} \\
%Simulate the replicator dynamics with $c=2, b=1$.  Start out with several
%sets of initial conditions ({\it i.e.} $(p,q)=(1,0), (0.2,0.8), (0,1)$). What
%happens? Does it depend on initial conditions? Does it agree with our Nash
%equilibrium from class?}}}
%\end{center}

\subsection{The Prisoner's Dilemma}
\noindent Recall for the prisoner's dilemma we had the payoff matrix
\[
A=\left( \begin{array}{cc}
-1 & -10 \\
0 & -7 \end{array} \right).
\]


\bigskip
\begin{center}
\colorbox{Goldenrod}{\parbox{15cm}{
{
\textbf{\large Exercise 2}\vspace{1mm} \\
Simulate the replicator dynamics for the Prisoner's Dilemma.  What
happens? Start out with several sets of initial conditions ({\it i.e.}
$(p,q)=(1,0), (0.2,0.8), (0,1)$). Does it agree with our Nash equilibrium from
class?}}}
\end{center}

\subsection{Snowdrift Game}
\noindent Recall the snowdrift game example from class,\footnote{Within the current framework where all players
play the same strategy, this example has no pure Nash equilibria. There is a
Nash equilibrium where one player plays the first strategy and the other plays
the second strategy (or \textit{vice versa}) but this is not considered in the
replicator dynamics framework.}
\[
A=\left(
\begin{array}{cc}
R& S \\
T& P
\end{array}
\right)
\]
with $T > R > S > P$.  

\bigskip
\begin{center}
\colorbox{Goldenrod}{\parbox{15cm}{
{
\textbf{\large Exercise 3}\vspace{1mm} \\
Simulate the replicator dynamics with $T = 4, R = 3, S = 2, P = 1.$  What happens? Start out with several
sets of initial conditions ({\it i.e.} $(p,q)=(1,0), (0.2,0.8), (0,1)$). Does
it agree with our result from class?}}}
\end{center}


\section{Analyzing these results}
\noindent We now want to understand more generally the results we found above. Looking
back at Eqs. \ref{eq:rep1} \& \ref{eq:rep2} we can see that they are redundant,
since we know that $p+q=1$. If we take the equation for $p$ and substitute
$q=1-p$ we get just a single equation for $p$,
 \begin{equation}
   \frac{dp}{dt}= p (1-p) \left( {\rm Payoff}_1-{\rm Payoff}_2\right). \label{eq:rep_sing1}
\end{equation}
In terms of our general payoff matrix, this is just
 \begin{equation}
   \frac{dp}{dt}= p (1-p) \left( (a_{11}-a_{12}-a_{21}+a_{22})p+a_{12}-a_{22}\right). \label{eq:rep_sing2}
\end{equation}
Now, the {\it equilibria} of the replicator dynamics are values of $p$ where
$$\frac{dp}{dt}=0.$$


\bigskip
\begin{center}
\colorbox{Goldenrod}{\parbox{15cm}{
{
\textbf{\large Exercise 4}\vspace{1mm} \\
Find the values of $p$ which can be equilibria.  Can you use this to
rationalize the results from the simulations above? {\it (Note that there will
be three possible equilibria, but only two of them always satisfy $0\le p\le
1$, as is required by the analysis.)}}}}
\end{center}

\section{The AM 50 Competition: The Unique Number Game}
\subsection{Introduction}
In the unique number game, there are $N$ players and each must choose an
integer between 0 and 99. The winner is the person who chose the smallest
number that was not chosen by anyone else. For example, if $N=5$ the
players might make the following choices:
\[
\begin{tabular}{c|c|c|c|c|c}
  Player & A & B & C & D & E \\
  \hline
  Choice & 0 & 5 & 2 & 3 & 0
\end{tabular}
\]

\noindent In this example, players A and E both chose zero so they invalidate each other.
Player C made the lowest remaining choice and is therefore the winner. In
certain rare cases, such as if B and D switched their choices to two, a game may
result in no winner.\\

\noindent  If the unique number game is played with human players, then choosing a number
requires some psychology about what the other players will pick. Lower numbers
are more desirable, but they are risky choices since there may be a greater
probability that others will also choose them. \\
%\textcolor{red}{While at the University of
%Cambridge, Chris played this game twice with $n=100$ undergraduate friends. The
%winning number in the first game was two, and the winning number in the second
%game was thirteen. In both games, about a quarter of the players chose zero.
%Interestingly, while the players were roughly evenly mixed in gender, almost
%everyone who chose zero was male.  Replace by what we saw in class.}\\

\noindent For many repeated games, we expect that a mixed-strategy Nash equilibrium may
exist, where different numbers are chosen with different probabilities.\\

\subsection{The game}
We are going to play many repeated unique number games in AM 50. Your job is to
write a Matlab function that implements a strategy in this repeated game. The
function should be called  \texttt{firstname\_lastname.m} and should accept five inputs
\texttt{N}, \texttt{history}, \texttt{lw}, \texttt{lc}, \texttt{y}.  These inputs will be
discussed in more detail later.  A simple strategy, given 
below, would be just to pick a fixed number:

\lstinputlisting{simple_strategy.m}

\noindent Alternatively, one could choose the numbers from 0 to 25 with
equal probability: 

\lstinputlisting{random_strategy.m}

\noindent Your function \texttt{firstname\_lastname.m}  should be submitted as part of the homework (due on Friday April 21st).\\

\noindent Once we have collected all of the programs we will run  (hopefully) one hundred sets of ten
million rounds each. The number of wins for each player will be counted across
all $10^9$ rounds. The winner will be rewarded with an amazing prize!

\subsection{Testing a strategy}
We have provided you with a program called \texttt{game\_test.m} that can be
used to test different strategies against each other.   It is currently set up to
simulate 10,000 rounds with $N=25$ players.  25 of the players use the
\texttt{random\_strategy.m} strategy, and the last player uses a very basic
strategy in \texttt{simple\_strategy.m}. Currently the player using the simple
strategy only wins a small fraction of approximately 0.38\% of the rounds, much less
than an equal share of \smash{$\frac{1}{25}=4\%$} of the rounds.

\bigskip
\begin{center}
\colorbox{Goldenrod}{\parbox{15cm}{
{
\textbf{\large Exercise 5}\vspace{1mm} \\
Modify \texttt{simple\_strategy.m} so that it wins consistently more than
6\% of the rounds.}}}
\end{center}

\subsection{More available information}
As in the two sample functions above, Your function  \texttt{firstname\_lastname.m} will be passed additional information about the
current state of the game and what has happened in previous rounds. If you
wish, you can make use of some of this information to devise a strategy. The
five inputs are as follows: \\


\begin{itemize}
  \item[\texttt{N}] The total number of players. \\
  \item[\texttt{history}] A list of length 100, containing the total number of
    choices of each number across all previous rounds that have been played so
    far.\\
  \item[\texttt{lw}] The winning number in the previous round. If there was no
    winner in the previous round this will be set to 100. In the very first
    round, before anything has been played, \texttt{lw} will be set to zero.\\
  \item[\texttt{lc}] A list of length \texttt{N} of all of the players' choices
    in the previous round. In the very first round, before anything has been
    played, all entries will be set to zero. Each player will occupy the same
    entry in the list throughout the game; for example, by repeatedly examining
    \texttt{lc[5]} you could see the entire history of one particular player.\\
  \item[\texttt{y}] Your position within the \texttt{lc} list. Hence
    \texttt{lc[y]} will be equal to your choice in the previous round.
\end{itemize}

\bigskip
\noindent Two functions that make use of some of this information are
provided for you to review:\\
\begin{itemize}
  \item[\texttt{prev.m}] If the previous winning choice was
    \texttt{lw}, then that may be a good number in general. This function
    returns \texttt{lw-1}, \texttt{lw}, or \texttt{lw+1} with equal
    probability.\\
 \item[\texttt{histo.m}] This function chooses the lowest number whose
    probability of occurrence in the previous games is less than
    \smash{$\frac{1}{N}$}, and adds a small random displacement.\\
 % \item[\texttt{means.m}] This function calculates the mean of the winning
   % choices in all previous games and chooses that, rounded to the nearest
    %integer.
\end{itemize}
\bigskip

\noindent The homework will have some additional exercises that may give you pointers on
devising a good strategy. After this, you can do the following: 


\bigskip
\begin{center}
\colorbox{Goldenrod}{\parbox{15cm}{
{
\textbf{\large Homework Question 1}\vspace{1mm} \\
Devise a strategy for the unique number game following the description and
rules above. Anything is acceptable, and you can even copy verbatim or modify
one of the basic examples discussed above. Submit it with the game theory
homework, with name \texttt{firstname\_lastname.m}.}}}
\end{center}
\bigskip

\subsection{Technicalities}
The following technicalities should not be relevant for any basic strategy but
may come into play if you want to try something more elaborate:\\
\begin{enumerate}
  \item We will review all submissions prior to playing the game. If we notice
    any functions that are liable to crash, we may contact you or make minor
    fixes.\\
  \item Designing strategy that involves cooperation beforehand with other
    students is allowed. For example, if you restrict your choice to odd
    numbers, and your friend restricts their choice to even numbers, then at
    least you know that you will never invalidate each other's choices.
    However, this a small advantage that may be outweighed by many other
    factors. Within the game, active communication between two strategies, such
    as by sharing data, is not allowed. The positions of each player within
    \texttt{lc} will be randomized, so there will be no straightforward way to
    identify your friend.\\
  \item If your function returns a value less than zero, it will be treated
    as choosing zero. If your function returns a value greater than 99, it will
    be treated as choosing 99.\\
  \item There will be a limit on how long your function can take. If a function
    (a) takes longer than 0.1~s more than 100 times in any set of rounds, (b)
    ever takes longer than 10~s, or (c) crashes, it will be disqualified and
    the game will be reset to not include this function.   The Matlab functions \texttt{tic} and \texttt{toc} can help you determine how long it takes your function to run.\\
%  \item The games will be run on machines with at least two gigabytes of
%    memory. In the unlikely event that the games run out memory, functions that
%    consume the most memory will be disqualified and the game will be reset.
  \item The teaching staff may disqualify any entry that is deemed to be
    subverting the rules or not in the spirit of the game.
\end{enumerate}








%
%\lstinputlisting{random_walkers.m}







\end{document}
